\documentclass[conference]{IEEEtran}
%%
\def\BibTeX{{\rm B\kern-.05em{\sc i\kern-.025em b}\kern-.08em
    T\kern-.1667em\lower.7ex\hbox{E}\kern-.125emX}}
%%
\usepackage{cite}
\usepackage{amsmath,amssymb,amsfonts}
\usepackage{url}
\usepackage{algorithmic}
\usepackage{graphicx}
\usepackage{textcomp}
\usepackage{xcolor}
\usepackage{lettrine}
\usepackage{listings}
\usepackage{booktabs}
\usepackage{hyperref}
\usepackage{pdfpages}

%% Reusable screenshot placeholder box (compiles without image files)
\newcommand{\screenshotbox}[2]{%
  \fbox{\parbox[c][#2][c]{0.95\linewidth}{\centering\textbf{Screenshot Placeholder}\\[0.4em]#1}}
}

%% Code listing style
\lstset{
  basicstyle=\ttfamily\small,
  breaklines=true,
  frame=single,
  backgroundcolor=\color{gray!10},
  keywordstyle=\color{blue},
  commentstyle=\color{green!60!black},
  stringstyle=\color{orange!80!black}
}

\begin{document}

%%% Uncomment the line below after obtaining your NCI Project Cover Sheet from Moodle
% \includepdf{nci-project-cover-sheet.pdf}

\title{SentriVault: A Cloud-Native Secure Secrets and Access Governance Platform}

\author{Ishan Verma\\
StudentId: [Your ID Here]\\
Cloud DevSecOps, MSc in Cloud Computing\\
National College of Ireland Dublin, IRELAND\\
Email: [your.email]@student.ncirl.ie\\
URL: \url{www.ncirl.ie}\\
Deployed Application URL: \url{[Your-Deployment-URL]}\\
GitHub Repository: \url{https://github.com/[your-repo]/DevOpSec}}

\maketitle

\begin{abstract}
Modern cloud apps store many secrets such as API keys, database passwords, and certificates. If these secrets leak, attackers can quickly access systems. This paper presents SentriVault, a cloud-native secrets management platform built with React, Node.js/Express, and PostgreSQL. The platform encrypts secrets before storage, controls access using role-based access control (RBAC), and records all key actions in an immutable audit log. SentriVault is deployed with GitHub Actions CI/CD to AWS Elastic Beanstalk (backend) and Amazon S3 (frontend). For static analysis, the project uses TypeScript strict mode, ESLint, and npm audit. During development, five security issues were found and fixed: invalid crypto input handling, incomplete cipher copy payload, revoked secrets still visible in list views, clipboard failures in non-HTTPS environments, and a missing audit action type. Testing confirmed that RBAC checks work correctly, revoked secrets are hidden, and there were no high or critical dependency vulnerabilities at submission time. The results show that strong secrets security can be achieved using standard open-source tools and cloud services.
\end{abstract}

\begin{IEEEkeywords}
secrets management, RBAC, AES-256-GCM, static code analysis, DevSecOps, CI/CD, audit logging
\end{IEEEkeywords}

%% ─────────────────────────────────────────────
\section{Introduction}
\label{sec:intro}

\lettrine[findent=2pt]{\fbox{\textbf{M}}}{ }odern cloud systems use many secrets: API keys, passwords, SSH keys, and certificates. In microservice projects, the number of secrets grows quickly. If even one secret is leaked, it can lead to a major breach. Recent reports show that stolen credentials remain one of the most common attack paths~\cite{verizon2023}. Still, many teams store secrets in plain text files, environment files, or source code. OWASP classifies these risks under \textit{Cryptographic Failures} (A02) and \textit{Security Misconfiguration} (A05)~\cite{owasp2021}.

SentriVault was built to show that a secure secrets platform can be created with common open-source tools and deployed on AWS without expensive enterprise vault products. The platform focuses on three security goals:

\begin{itemize}
  \item \textbf{Confidentiality:} Secret values are encrypted before persistence using symmetric authenticated encryption; plaintext is never stored and never surfaced through list endpoints.
  \item \textbf{Least-Privilege Access:} Every API route is gated by a unique RBAC permission key evaluated at runtime; even valid tokens lacking the required permission receive HTTP 403.
  \item \textbf{Accountability:} Every access event is logged to an immutable, append-only audit trail recording the actor, target, action, success state, source IP, and timestamp.
\end{itemize}

This report covers the design, implementation, testing, and deployment of SentriVault. The stack uses React + TypeScript for frontend, Express + Node.js for backend, PostgreSQL for data, and GitHub Actions for CI/CD. The design follows defence-in-depth: encrypt data, enforce RBAC at runtime, validate crypto input, and log important actions.

Section~\ref{sec:cicd} details the CI/CD pipeline architecture, deployment workflow, and integration with GitHub Actions. Section~\ref{sec:sca} presents the static code analysis approach, identifies five key vulnerabilities discovered and fixed, and documents the findings from TypeScript, ESLint, and dependency scanning. Section~\ref{sec:conclusions} interprets results and recommends future hardening priorities.

%% ─────────────────────────────────────────────
\section{Project Requirements}
\label{sec:requirements}

\lettrine[findent=2pt]{\fbox{\textbf{T}}}{ }his section explains the required project stack and how each layer works together in SentriVault.

\subsection{Required Project Stack}

The implementation uses the following stack:

\begin{itemize}
  \item Frontend: React 18, TypeScript 5, Vite 5, React Router v6
  \item Backend: Node.js 20, Express 5, TypeScript 5, Zod validation
  \item Database: PostgreSQL 16 with role-based schema and audit tables
  \item Security: JWT authentication, bcrypt password hashing, RBAC middleware, secret encryption utilities
  \item DevOps: Docker Compose for local database, GitHub Actions for CI/CD, AWS Elastic Beanstalk for backend deployment, Amazon S3 + CloudFront for frontend deployment
  \item Quality and security checks: TypeScript strict mode, ESLint, npm audit (and optional SonarQube/SonarCloud)
\end{itemize}

\subsection{How the Stack Works Together}

The user opens the React web app and signs in. The app sends API requests to the Express backend. The backend validates the JWT, checks RBAC permissions, validates request body with Zod, and then reads or writes data in PostgreSQL. For secrets, the backend encrypts values before saving and decrypts only on valid requests. Every important action is written to the audit log table.

In development, Docker Compose runs PostgreSQL locally while frontend and backend run with Node tooling. In production, GitHub Actions builds both apps and deploys backend and frontend to AWS services. This flow keeps development simple and deployment repeatable.

\subsection{Project Code Directory Evidence}

\begin{figure}[ht]
\begin{center}
\screenshotbox{Insert screenshot of project directory structure here.\\Capture folders such as backend, frontend, docs, infra, and .github/workflows from your IDE file explorer.}{5.0cm}
\caption{Project code directory screenshot placeholder.\label{fig:project-directory-shot}}
\end{center}
\end{figure}

Recommended screenshot content for Figure~\ref{fig:project-directory-shot}:

\begin{itemize}
  \item Root folders: \texttt{backend}, \texttt{frontend}, \texttt{docs}, \texttt{infra}, \texttt{.github}
  \item Backend source modules: \texttt{auth}, \texttt{secrets}, \texttt{users}, \texttt{roles}, \texttt{audit}
  \item Frontend pages and context folders
  \item CI/CD workflow file under \texttt{.github/workflows}
\end{itemize}

%% ─────────────────────────────────────────────
\section{Continuous Integration, Continuous Delivery and Deployment}
\label{sec:cicd}

\lettrine[findent=2pt]{\fbox{\textbf{T}}}{ }his section documents the CI/CD pipeline, including its architecture, the workflow of code changes through the system, and the tooling employed.

\subsection{Architecture Overview}

SentriVault is a three-tier application: a React 18 single-page application (SPA) built with Vite communicates over HTTPS with a REST API implemented in Node.js and Express, which persists all state in PostgreSQL. Figure~\ref{fig:arch} illustrates the deployment topology.

\begin{figure}[ht]
\begin{center}
\includegraphics[width=\linewidth]{architecture-placeholder.png}
\caption{SentriVault three-tier architecture. React SPA hosted on S3 + CloudFront communicates with Express API on Elastic Beanstalk backed by PostgreSQL.\label{fig:arch}}
\end{center}
\end{figure}

In production, the frontend is served as a static Vite bundle from Amazon S3 with CloudFront CDN caching; the backend runs in AWS Elastic Beanstalk behind an Application Load Balancer. For local development, Docker Compose instantiates a PostgreSQL container on the developer's machine. Table~\ref{table:stack} lists the core dependencies.

\begin{figure}[ht]
\begin{center}
\screenshotbox{Insert your deployed website home/dashboard screenshot here.\\Example file name: website-dashboard.png}{4.2cm}
\caption{Website screenshot placeholder (UI evidence).\label{fig:website-shot}}
\end{center}
\end{figure}

\begin{table}[h]
\centering
\caption{Technology Stack\label{table:stack}}
\renewcommand{\arraystretch}{1.05}
\begin{tabular}{||l l||} 
\hline
\textbf{Component} & \textbf{Technology} \\[0.5ex] 
\hline\hline
Frontend Framework & React 18 + TypeScript 5 \\
Build Tool & Vite 5 \\
Routing & React Router v6 \\
API Framework & Express 5 \\
Backend Language & Node.js 20 + TypeScript 5 \\
Request Validation & Zod 3 \\
Database & PostgreSQL 16 \\
Authentication & JSON Web Tokens (JWT) \\
Password Hashing & bcrypt (cost 12) \\
Database Driver & node-postgres (pg) \\
Containerization & Docker Compose 3.9 \\
CI/CD Orchestration & GitHub Actions \\
Infrastructure & AWS (Elastic Beanstalk, S3, CloudFront, RDS) \\
\hline
\end{tabular}
\end{table}

\subsection{CI/CD Pipeline Design}

The GitHub Actions workflow (stored in \texttt{.github/workflows/ci-cd.yml}) implements a three-stage pipeline: Continuous Integration (build and test), Continuous Delivery (backend to Elastic Beanstalk), and Continuous Delivery (frontend to S3).

\begin{figure}[ht]
\begin{center}
\screenshotbox{Insert your full CI/CD workflow screenshot here (GitHub Actions jobs and stages).\\Include CI and both CD jobs in one view.}{4.8cm}
\caption{Pipeline overview screenshot placeholder.\label{fig:pipeline-overview}}
\end{center}
\end{figure}

\subsubsection{Continuous Integration}

On every push to the \texttt{main} branch, the CI job:

\begin{enumerate}
  \item Checks out the repository
  \item Sets up Node.js 20
  \item Installs backend dependencies via \texttt{npm ci}
  \item Compiles backend TypeScript with \texttt{npm run build}
  \item Installs frontend dependencies via \texttt{npm ci}
  \item Compiles frontend TypeScript with \texttt{npm run build}
\end{enumerate}

If any step fails, the job fails immediately and no deployment proceeds. TypeScript \texttt{strict: true} mode acts as an enforced static type check: any type error aborts the build. This constraint prevents entire categories of bugs (implicit \texttt{any}, unsafe property access, missing null checks) from reaching production.

\subsubsection{Continuous Deployment — Backend}

On successful CI completion, the CD Backend job:

\begin{enumerate}
  \item Zips the built backend application (excluding \texttt{node\_modules/})
  \item Generates a versioned deploy label: \texttt{devopsec-backend-\$\{git\ sha\}-r\$\{run\_id\}}
  \item Uploads the ZIP to an S3 bucket
  \item Invokes the AWS Elastic Beanstalk API to deploy the new version
\end{enumerate}

The backend deployment uses environment variables injected from GitHub Actions secrets at runtime: \texttt{ENCRYPTION\_KEY}, \texttt{DATABASE\_URL}, \texttt{JWT\_SECRET}, \texttt{AWS\_ACCESS\_KEY\_ID}, \texttt{AWS\_SECRET\_ACCESS\_KEY}. These credentials never appear in repository files.

\subsubsection{Continuous Deployment — Frontend}

The CD Frontend job:

\begin{enumerate}
  \item Builds the Vite production bundle via \texttt{npm run build}
  \item Syncs the \texttt{dist/} directory to S3 using \texttt{aws s3 sync}
  \item Invalidates the CloudFront distribution cache
\end{enumerate}

The deployed frontend is served from \texttt{https://devopsec.example.com} (substitute your actual domain). All API requests go to the backend at \texttt{https://api.devopsec.example.com/api/v1}.

\subsubsection{Production Result Evidence}

To clearly show production success, include both pipeline and live-system evidence:

\begin{figure}[ht]
\begin{center}
\screenshotbox{Insert screenshot of live production application in browser.\\Show deployed URL and one core page (dashboard/secrets).}{4.2cm}
\caption{Production application screenshot placeholder.\label{fig:prod-app-shot}}
\end{center}
\end{figure}

\begin{figure}[ht]
\begin{center}
\screenshotbox{Insert screenshot of production backend status.\\Example: Elastic Beanstalk health green, ALB target healthy, or API health endpoint response.}{4.2cm}
\caption{Production backend status screenshot placeholder.\label{fig:prod-backend-shot}}
\end{center}
\end{figure}

\subsection{Deployment Function — Code Change Flow}

A developer making changes to the codebase follows this workflow:

\begin{enumerate}
  \item Creates a feature branch from \texttt{main}
  \item Modifies source files (e.g., \texttt{backend/src/modules/secrets/routes.ts})
  \item Commits and pushes the branch to GitHub: \texttt{git push origin feature/my-change}
  \item Opens a pull request against \texttt{main}
  \item The CI workflow triggers automatically on the PR, building both backend and frontend
  \item Upon approval and merge, the code merges to \texttt{main}
  \item The CI workflow runs again on \texttt{main}; on success, the CD jobs deploy to Elastic Beanstalk and S3
\end{enumerate}

Example: A developer fixes a bug in the decrypt route by modifying \texttt{backend/src/modules/crypto/routes.ts} (adding explicit IV length validation). The change is committed and pushed. GitHub Actions automatically compiles the TypeScript, runs the build, and if successful, deploys the new backend to Elastic Beanstalk. Within seconds, all incoming decrypt requests use the patched code.

\begin{figure}[ht]
\begin{center}
\screenshotbox{Insert a screenshot of one successful pipeline run after a code change.\\Show commit ID, green checks, and deployment steps.}{4.8cm}
\caption{Pipeline run-in-action screenshot placeholder.\label{fig:pipeline-run}}
\end{center}
\end{figure}

\begin{figure}[ht]
\begin{center}
\screenshotbox{Insert source code change screenshot from your IDE (before/after fix).\\Good example: decrypt IV validation or revoked-secret filter.}{4.5cm}
\caption{Code-change evidence screenshot placeholder.\label{fig:code-change-shot}}
\end{center}
\end{figure}

\subsection{Database and Local Development}

Local development begins with:

\begin{lstlisting}[language=bash]
docker-compose up -d
cd backend && npm install && npm run build && npm run dev
cd frontend && npm install && npm run dev
\end{lstlisting}

This launches a PostgreSQL container on \texttt{localhost:5432}, the Express server on \texttt{localhost:4000}, and the Vite dev server on \texttt{localhost:5173}. The frontend's \texttt{.env.local} sets \texttt{VITE\_API\_BASE\_URL=http://localhost:4000/api/v1}.

All database schema migrations are version-controlled in \texttt{backend/prisma-or-sql/} and applied manually or via deployment scripts. The schema implements the RBAC lattice (roles, permissions, user-role mappings), secret storage with versioning, and append-only audit logging.

%% ─────────────────────────────────────────────
\section{Static Code Analysis and Security Vulnerability Analysis}
\label{sec:sca}

\lettrine[findent=2pt]{\fbox{\textbf{T}}}{ }his section briefly documents tools used, key findings, and fixes.

\subsection{Tools Used}

We used three core checks in CI and local verification:

\begin{itemize}
  \item \textbf{TypeScript strict mode} to catch type-safety issues before deployment.
  \item \textbf{ESLint} to detect unsafe code patterns and code-quality issues.
  \item \textbf{npm audit} to detect known dependency vulnerabilities.
\end{itemize}

\begin{figure}[ht]
\begin{center}
\screenshotbox{Insert SonarQube/SonarCloud summary screenshot.\\Keep only project overview, issue count, and quality gate.}{3.4cm}
\caption{SonarQube/SonarCloud result placeholder.\label{fig:sonar-shot}}
\end{center}
\end{figure}

\subsection{Key Findings and Fixes (Summary)}

\begin{table}[h]
\centering
\caption{Static analysis and security fixes summary\label{table:vulns}}
\renewcommand{\arraystretch}{1.1}
\begin{tabular}{||p{1.1cm} p{2.0cm} p{2.8cm}||}
\hline
\textbf{Type} & \textbf{Finding} & \textbf{Fix} \\[0.5ex]
\hline\hline
Crypto & Invalid IV handling in decrypt path & Added explicit IV/auth-tag validation before decrypt \\
\hline
Crypto & Incomplete cipher copy payload & Copied full JSON payload: value, iv, auth\_tag, algorithm \\
\hline
Access & Revoked secrets still listed & Added query filter: status not equal to revoked \\
\hline
Runtime & Clipboard API failed in non-HTTPS & Added fallback copy method with \texttt{execCommand} \\
\hline
Type safety & Missing audit action type & Added \texttt{update\_secret} to \texttt{AuditAction} union \\
\hline
\end{tabular}
\end{table}

\subsection{Result Snapshot}

Final checks showed one type issue fixed, no high-priority ESLint errors, and no high/critical npm audit vulnerabilities at submission time.

\begin{figure}[ht]
\begin{center}
\screenshotbox{Insert compact static-analysis evidence screenshot.\\Example: terminal output of tsc + eslint + npm audit summary.}{3.4cm}
\caption{Static analysis output placeholder.\label{fig:static-summary-shot}}
\end{center}
\end{figure}

%% ─────────────────────────────────────────────
\section{Conclusions}
\label{sec:conclusions}

\lettrine[findent=2pt]{\fbox{\textbf{T}}}{ }his project shows that you can build a secure secrets platform using common open-source tools and AWS services. SentriVault uses three strong controls: encryption for confidentiality, RBAC for access control, and audit logs for traceability. CI/CD and static analysis helped catch issues early and keep deployments safe.

\subsection{Key Findings}

\begin{enumerate}
  \item \textbf{Static analysis catches real defects:} TypeScript strict mode identified the missing \texttt{update\_secret} audit action value, a bug that dynamic testing might have missed until the specific code path was exercised.
  \item \textbf{Defence-in-depth is practical:} Layering encryption, RBAC, auditing, and CI/CD validation produced measurably more secure secrets handling than ad-hoc environment variables.
  \item \textbf{Zero high-severity CVEs achieved:} Dependency scanning and ESLint integration blocked vulnerable packages from production.
  \item \textbf{RBAC runtime enforcement scales:} Live database permission checks (rather than JWT claims) enable immediate revocation without token rotation.
\end{enumerate}

\subsection{Reflection and Lessons Learned}

If this project is implemented again, these improvements should be done first:

\begin{itemize}
  \item \textbf{Key management:} Use AWS KMS or HashiCorp Vault for hardware-backed key storage rather than environment variables. This eliminates the single-server key-compromise risk.
  \item \textbf{PBKDF2 key derivation:} Replace SHA-256 hashing with PBKDF2 or \texttt{scrypt} and a salt to harden against brute-force attacks on key material.
  \item \textbf{Test coverage gates:} Add minimum coverage thresholds (e.g., 80\%) enforced in CI. Current integration tests lack full workflow coverage (secret creation $\to$ rotation $\to$ revocation).
  \item \textbf{Rate limiting:} Implement rate limiting on authentication endpoints to prevent credential-stuffing attacks.
  \item \textbf{Refresh token rotation:} Complete implementation of JWT refresh-token rotation so that compromised tokens have bounded lifetime.
  \item \textbf{SonarQube/SonarCloud integration:} Add automated code-quality scanning to the CI pipeline to surface maintainability and security-pattern issues beyond type checking and linting.
  \item \textbf{Runtime secrets monitoring:} Log and alert on unusual access patterns (e.g., a user accessing thousands of secrets in one session) to detect account compromise or lateral movement.
\end{itemize}

\subsection{Future Directions}

\begin{enumerate}
  \item \textbf{HashiCorp Vault integration:} Extend SentriVault to act as a management layer on top of Vault, enabling centralized key rotation and policy management.
  \item \textbf{Hardware security modules (HSM):} Support AWS CloudHSM and other HSM providers for key material that never leaves the hardware.
  \item \textbf{Secrets rotation automation:} Implement scheduled automatic rotation of database passwords, API keys, and certificates, with automatic downstream updates to dependent services.
  \item \textbf{ML-powered anomaly detection:} Analyze audit logs with machine learning to identify suspicious access patterns and raise alerts.
  \item \textbf{Multi-region replication:} Extend the architecture to replicate secrets and audit logs across AWS regions for disaster recovery and compliance.
\end{enumerate}

In summary, SentriVault provides a solid foundation for managing secrets at scale with verifiable security controls. The integration of static analysis into the deployment pipeline, combined with runtime access control and immutable auditing, offers a significantly stronger posture than traditional ad-hoc secret storage practices.

\subsection{Screenshot Checklist (for final submission)}

Before submission, replace all placeholders with your own screenshots and confirm these are included:

\begin{itemize}
  \item Website UI (landing page and dashboard)
  \item CI/CD workflow overview (all jobs)
  \item One pipeline run after a real code change
  \item Code change in IDE (security fix)
  \item SonarQube/SonarCloud project dashboard
  \item TypeScript build output
  \item ESLint report/output
  \item npm audit report for backend and frontend
\end{itemize}

%% ─────────────────────────────────────────────

\bibliography{cdos-biblio}
\bibliographystyle{IEEEtran}

\end{document}
